\begin{table}[H]\centering\small
\caption{更强基线：同初末6000 kWh的跨日贪心与主策略（万元）}
\begin{tabular}{lrrrrr}\toprule
机制&跨日贪心&主策略&节省&降幅&7日块95\%区间\\\midrule
问题二 & 1404.66 & 1384.42 & 20.24 & 1.44\% & [14.94,25.83] \\
问题三 & 1369.31 & 1346.18 & 23.13 & 1.69\% & [17.80,27.61] \\
问题四日前 & 1480.75 & 1465.46 & 15.29 & 1.03\% & [10.50,19.66] \\
问题四日内 & 1448.10 & 1425.21 & 22.89 & 1.58\% & [17.69,27.09] \\
\bottomrule\end{tabular}\end{table}
跨日贪心采用场景共同合同与即时电量平衡反馈，保留中间午夜连续性，仅在真正年末恢复6000 kWh。主策略相对该基线节省15.29至23.13万元，降幅1.03\%至1.69\%，且四组七日块重采样下界均为正。这是在同样允许跨日的条件下，合同规划与库存反馈整体框架相对贪心执行的差额，不能与主表每日闭合基线的节省相加，也不能把全部差额归因于单独一层。
